\chapter{Results and Discussion of Ansieyes}

We evaluated Ansieyes across its four core capabilities: issue classification quality, duplicate detection accuracy, prompt injection defence, and token efficiency of the two-pass architecture. This chapter presents our findings along with a discussion of what worked well and where the system has room for improvement.

\section{Experimental Setup}

All experiments used Gemini 2.0 Flash as the primary model. The test repository was the Ansieyes codebase itself (roughly 5,000 lines across 25+ source files), supplemented by a curated set of 33 real and synthetic issues. We ran the \gls{ci} pipeline across Python 3.11, 3.12, and 3.13 to confirm cross-version compatibility.

\section{Issue Classification}

We tested the system on 33 issues with manually verified ground-truth labels. Table \ref{tab:classification} breaks down the results by issue type.

\begin{table}[htb]
\centering
\renewcommand{\arraystretch}{1.3}
\setlength{\arrayrulewidth}{0.5mm}
\fontsize{10}{12}\selectfont
\caption{\textbf{Issue classification accuracy by type}}
\vspace{2mm}
\label{tab:classification}
\begin{tabular}{|p{3.5cm}|c|c|c|}
\hline
\textbf{Type} & \textbf{Total} & \textbf{Correct} & \textbf{Accuracy} \\
\hline
Bug & 15 & 14 & 93.3\% \\
\hline
Enhancement & 10 & 9 & 90.0\% \\
\hline
Feature Request & 8 & 7 & 87.5\% \\
\hline
\textbf{Overall} & \textbf{33} & \textbf{30} & \textbf{90.9\%} \\
\hline
\end{tabular}
\end{table}

Bug classification had the highest accuracy, which makes sense---bug reports tend to contain more distinctive language (error messages, stack traces, ``crash'' or ``fails'') that gives the model strong signals. Feature requests were occasionally confused with enhancements, particularly when the issue described extending an existing capability rather than building something new. The single misclassified bug was a performance issue that the model labelled as an enhancement, a reasonable but incorrect judgement.

\subsection{Confidence Scores}
Table \ref{tab:confidence} shows how the system's self-reported confidence scores distributed across analyses.

\begin{table}[htb]
\centering
\renewcommand{\arraystretch}{1.3}
\setlength{\arrayrulewidth}{0.5mm}
\fontsize{10}{12}\selectfont
\caption{\textbf{Distribution of confidence scores}}
\vspace{2mm}
\label{tab:confidence}
\begin{tabular}{|c|c|c|}
\hline
\textbf{Range} & \textbf{Count} & \textbf{Percentage} \\
\hline
0.8 -- 1.0 & 18 & 54.5\% \\
\hline
0.6 -- 0.79 & 10 & 30.3\% \\
\hline
0.4 -- 0.59 & 3 & 9.1\% \\
\hline
0.3 -- 0.39 & 2 & 6.1\% \\
\hline
\end{tabular}
\end{table}

We found that analyses with confidence above 0.7 almost always produced correct classifications and useful root cause descriptions. The two lowest-confidence results came from issues with very brief descriptions (one or two sentences), where the model simply did not have enough information to work with.

\subsection{Root Cause and Solution Quality}
The system identified at least one relevant code location for 28 out of 33 issues (84.8\%). Manual review confirmed that the primary cause was accurate in 27 cases (82\%). The five misses were split between issues involving configuration files (which Repomix sometimes excluded) and issues that spanned multiple repositories.

At least one proposed solution was present in 97\% of analyses. The retry mechanism caught the remaining 3\% on the second attempt, so after retries, every issue received at least one solution proposal.

\section{Duplicate Detection}

We constructed a test set of 20 known duplicate pairs and 30 non-duplicate pairs, then ran both detectors on them.

\begin{table}[htb]
\centering
\renewcommand{\arraystretch}{1.3}
\setlength{\arrayrulewidth}{0.5mm}
\fontsize{10}{12}\selectfont
\caption{\textbf{Duplicate detection: Gemini vs. Cosine similarity}}
\vspace{2mm}
\label{tab:dupcomp}
\begin{tabular}{|p{4cm}|c|c|}
\hline
\textbf{Metric} & \textbf{Gemini} & \textbf{Cosine} \\
\hline
Precision & 0.89 & 0.78 \\
\hline
Recall & 0.85 & 0.70 \\
\hline
F1 Score & 0.87 & 0.74 \\
\hline
Avg. Detection Time & 2.3 s & $<$ 0.1 s \\
\hline
\end{tabular}
\end{table}

The Gemini-based detector handled semantic duplicates well---it correctly matched issues that described the same bug using entirely different wording. Its main source of error was overly cautious behaviour: our conservative prompt instruction (``prefer not marking as duplicate when uncertain'') caused it to miss a few genuine duplicates with low surface similarity.

The cosine detector's lower recall was expected. It struggles with paraphrased descriptions because \gls{tfidf} only captures word-level overlap. However, its near-instant speed makes it useful as a pre-filter in high-volume scenarios: run cosine first to catch the obvious duplicates, then use Gemini for the remaining cases.

\section{Prompt Injection Detection}

We tested the security module against 50 injection attempts (spread across all seven categories) and 50 benign issue descriptions.

\begin{table}[htb]
\centering
\renewcommand{\arraystretch}{1.3}
\setlength{\arrayrulewidth}{0.5mm}
\fontsize{10}{12}\selectfont
\caption{\textbf{Injection detection rates by category}}
\vspace{2mm}
\label{tab:injection}
\begin{tabular}{|p{3.5cm}|c|c|c|}
\hline
\textbf{Category} & \textbf{Tested} & \textbf{Caught} & \textbf{Rate} \\
\hline
Role Manipulation & 8 & 8 & 100\% \\
\hline
System Prompts & 7 & 7 & 100\% \\
\hline
Instruction Bypass & 8 & 7 & 87.5\% \\
\hline
File Manipulation & 6 & 6 & 100\% \\
\hline
Code Injection & 7 & 7 & 100\% \\
\hline
Data Extraction & 7 & 6 & 85.7\% \\
\hline
Prompt Leakage & 7 & 7 & 100\% \\
\hline
\textbf{Overall} & \textbf{50} & \textbf{48} & \textbf{96\%} \\
\hline
\end{tabular}
\end{table}

The two missed injections used subtle, context-dependent phrasing that avoided all seven regex categories and fell below the heuristic thresholds. On the benign side, we saw 2 false positives out of 50 (4\%). Both were security-focused issues that naturally contained words like ``bypass'' and ``injection'' in their technical descriptions. This is a known trade-off with keyword-based detection---legitimate security discussions can look suspicious to a pattern matcher.

\section{Two-Pass Token Efficiency}

Table \ref{tab:tokens} compares token usage between the single-pass and two-pass approaches.

\begin{table}[htb]
\centering
\renewcommand{\arraystretch}{1.3}
\setlength{\arrayrulewidth}{0.5mm}
\fontsize{10}{12}\selectfont
\caption{\textbf{Token consumption: single-pass vs. two-pass}}
\vspace{2mm}
\label{tab:tokens}
\begin{tabular}{|p{3cm}|c|c|c|}
\hline
\textbf{Repo Size} & \textbf{Single-Pass} & \textbf{Two-Pass} & \textbf{Reduction} \\
\hline
$<$5K LoC & 15,000 & 12,000 & 20\% \\
\hline
5--20K LoC & 65,000 & 28,000 & 57\% \\
\hline
20--100K LoC & 250,000 & 75,000 & 70\% \\
\hline
$>$100K LoC & Exceeds limit & 120,000 & --- \\
\hline
\end{tabular}
\end{table}

For small repositories, the two-pass architecture adds overhead (the Librarian pass itself costs tokens) without much benefit. The sweet spot is repositories in the 5K--100K line range, where the Librarian reliably prunes irrelevant files and the Surgeon receives a tightly focused context. For repositories exceeding 100K lines, the single-pass approach is simply not feasible, making the two-pass the only option.

\section{Execution Time}

Table \ref{tab:timing} summarises the wall-clock time for each operation.

\begin{table}[htb]
\centering
\renewcommand{\arraystretch}{1.3}
\setlength{\arrayrulewidth}{0.5mm}
\fontsize{10}{12}\selectfont
\caption{\textbf{Execution time by operation}}
\vspace{2mm}
\label{tab:timing}
\begin{tabular}{|p{5cm}|c|}
\hline
\textbf{Operation} & \textbf{Time} \\
\hline
Security check & $<$ 1 s \\
\hline
Cosine duplicate detection & $<$ 0.1 s \\
\hline
Gemini duplicate detection & 1--5 s \\
\hline
Single issue analysis & 10--30 s \\
\hline
Librarian pass & 5--15 s \\
\hline
Full two-pass analysis & 15--45 s \\
\hline
\end{tabular}
\end{table}

The dominant factor in total execution time is the Gemini \gls{api} latency. The security check and cosine duplicate detection add negligible overhead. Even the full two-pass analysis completes within a minute, which is well within the timeout limits of GitHub Actions workflows.

\section{CI Pipeline}

All unit tests pass across Python 3.11, 3.12, and 3.13. Code formatting (Black), import ordering (isort), and linting (Flake8) checks also pass consistently. The parallel test matrix catches version-specific issues early---we encountered one such issue with a Pydantic behaviour change in Python 3.13 and fixed it before it reached production.

The experimental results confirm that Ansieyes meets its design objectives. The system provides accurate, context-aware analysis for the majority of issues, detects most adversarial inputs, and scales to large codebases through the two-pass architecture. The next chapter summarises our conclusions and outlines future directions.

\section{Summary}

This chapter presented the experimental evaluation of Ansieyes across its four core capabilities. Issue classification achieved 90.9\% accuracy, with root cause analysis correctly identifying relevant code locations in 85\% of cases. Duplicate detection reached an F1 of 0.87 with the Gemini-based approach, while the cosine similarity fallback provided near-instant pre-screening. The prompt injection defence caught 96\% of attacks with a 4\% false positive rate. The two-pass Librarian--Surgeon architecture reduced token consumption by 57--70\% for medium to large repositories, enabling analysis of codebases that exceed the single-pass context limit.
